% !TITLE 破阵子·为陈同甫赋壮词以寄之
% !AUTHOR 辛弃疾
% !DYNASTY 两宋
% !GENRE 词
% !ORDER 17

\begin{poem}
醉里挑灯\textsuperscript{1}看剑，梦回吹角连营\textsuperscript{2}。
八百里分麾下炙\textsuperscript{3}，五十弦翻塞外声\textsuperscript{4}，沙场秋点兵\textsuperscript{5}。

马作的卢飞快\textsuperscript{6}，弓如霹雳弦惊\textsuperscript{7}。
了却君王天下事\textsuperscript{8}，赢得生前身后名，可怜白发生！
\end{poem}

\begin{annotations}
\notetext{1}{挑灯：拨亮灯芯。挑（tiǎo），把灯芯往上拨使灯火更亮。酒醉了还要拨亮灯来看剑——一个渴望上战场的将军形象跃然纸上。}
\notetext{2}{吹角连营：连绵军营中响起的号角声。角，古代军中的号角。}
\notetext{3}{八百里分麾下炙：把烤牛分给部下享用。“八百里”指牛，典出晋代王恺的牛名“八百里驳”；麾下，部下；炙（zhì），烤肉。}
\notetext{4}{五十弦翻塞外声：五十弦，指瑟（传说瑟有五十弦）；翻，弹奏；塞外声，边塞的军歌。军乐奏响边塞的雄壮之音。}
\notetext{5}{沙场秋点兵：在秋天的战场上检阅军队。点兵，阅兵。}
\notetext{6}{的卢：一种名马的名字。三国时刘备骑的卢马，在危急时一跃三丈过檀溪脱险。“马作的卢飞快”——战马像的卢一样飞驰。}
\notetext{7}{霹雳：惊雷。弓弦声像霹雳一样震耳。用雷声比弓弦声，写战斗的激烈。}
\notetext{8}{天下事：恢复中原、统一国家的大事。了却君王的心事，赢得身前身后的功名——但这一切只是梦，现实是“可怜白发生”，自己已经老了。}
\end{annotations}

\begin{appreciation}
《破阵子·为陈同甫赋壮词以寄之》是辛弃疾写给老朋友陈亮的词，也是他一生里最有名的壮词之一——写的是一场梦。

开头就耐人寻味："醉里挑灯看剑"——喝醉了，还要把灯芯挑亮，把剑抽出来看一看。醉醺醺的将军半夜看剑，看的是剑，想的是战场。这一看，就"梦回吹角连营"了：连片的军营，号角连天。八百里分麾下炙，五十弦翻塞外声，沙场秋点兵——把烤牛肉分给部下，军乐奏着边塞的调子，秋天的大校场上，听他一个人检阅。"八百里"是晋代一头名牛的名字，"五十弦"是瑟的典故，辛弃疾用起典来不在乎生僻，他要的是那股劲儿。

下阕打得正酣：马像的卢一样飞快，弓弦响得像炸雷。了却君王天下事，赢得生前身后名——替君王完成收复中原的大业，挣一个流传后世的名声。到这里，梦做到了最高处。

然后，一句"可怜白发生！"当头砸下来。原来前面写的全是想象。辛弃疾二十岁出头就带兵杀进金营，活捉叛徒，率众南归，以为到了朝廷就能挥师北伐；结果四十多年，他上书的北伐方略一条也没被采纳，大半时间闲居在上饶、铅山，种地看云。他和陈亮，两个满脑子北伐的人，只能借一阕词，来一场醉里的阅兵。

还记得《无衣》里那些"与子同袍"的战士吗？三千年前，他们修好戈矛就能上阵，同仇、偕作、偕行，一声比一声急。辛弃疾却连战场都上不了，刀在鞘里生锈，人先白了头。李煜也填过一阕《破阵子》，写的是亡国；辛弃疾这阕《破阵子》，写的是上不了战场的将军梦。同一个词牌，一个痛在失去，一个痛在不得。

"可怜白发生"是给所有人的提醒：想做的时候没做，等想做了，头发先白了。你还小，正是最来得及的年纪。想学的东西、想做的事，别攒着。哥哥不是催你干大事，只是希望你将来的"后悔清单"能短一点。
\end{appreciation}
